\subsubsection{Classical running of neutrino masses from six dimensions --- arXiv:hep-ph/0511001}

\textbf{Authors:} E. Dudas, C. Grojean, S. K. Vempati

\paragraph{主な主張}
$T^2/\mathbb{Z}_2$上の6次元bulk Weylフェルミオンとbrane上のactiveニュートリノを局在Dirac質量で結ぶ模型では、質量スペクトルにbraneの厚さゼロ極限に由来する対数的UV依存が生じ、局在結合の古典的なくりこみ群によるrunningとして解釈できることを示した\cite{Dudas:2005ClassicalRunning}。

\paragraph{新規性}
スカラー場で知られていたcodimension 2の局在結合の古典的runningをニュートリノ質量生成へ拡張し、KK質量行列の対角化と境界条件から同じ固有値方程式を導いた。

\paragraph{位置付け}
2余剰次元のbulkニュートリノ模型における局在相互作用の正則化・くりこみを扱い、KKスペクトルと振動のUV依存を検討する基礎となる。
